\chapter{Glosario}\label{app:glosario}

Los términos que más se repiten en el informe, explicados en pocas palabras. Están agrupados por tema y, dentro de
cada grupo, en orden alfabético.

\section*{Sensores y señal}
\begin{description}[style=sameline,leftmargin=1.5em,labelindent=0pt]
  \item[DC bias (polarización)] Tensión continua sobre la que ``flota'' la señal. En un sensor IEPE es la tensión que
  aparece en el cable cuando el sensor está en reposo (unos \SIrange{8}{14}{\volt}); la señal acústica es una pequeña
  variación encima de ella.
  \item[Ganancia de ruido] Factor por el que un amplificador multiplica su propio ruido y su offset de entrada. En el
  sumador de esta placa es $1 + R_f/R_\text{in} \approx \num{1.1}$.
  \item[IEPE] \emph{Integrated Electronics Piezo-Electric}. Sensor con un preamplificador interno que se alimenta con
  una corriente constante por el mismo coaxial por el que entrega la señal. Necesita una fuente de corriente (aquí,
  un LT3092 por canal).
  \item[Lazo de DC] Corrección automática del nivel de continua: el microcontrolador mide \net{DC_MON}, calcula el
  error y ajusta el DAC, que mueve \net{VSHIFT} hasta que la salida queda en el valor deseado.
  \item[Sumador inversor] Configuración de op-amp que suma varias entradas con signo cambiado. Aquí resta la
  polarización del micrófono y deja solo la señal, centrada en un nivel que fija \net{VSHIFT}.
\end{description}

\section*{Alimentación}
\begin{description}[style=sameline,leftmargin=1.5em,labelindent=0pt]
  \item[Buck] Convertidor conmutado que baja la tensión (de \SI{24}{\volt} a \SI{3.3}{\volt} o \SI{5}{\volt}) prendiendo y
  apagando un transistor muy rápido y filtrando con una bobina y condensadores. Pierde mucha menos energía que un
  regulador lineal.
  \item[CCM] \emph{Continuous Conduction Mode}. La corriente de la bobina nunca llega a cero dentro de un ciclo. En FPWM
  el LMR36015 trabaja siempre así.
  \item[Ciclo útil ($D$)] Fracción del periodo en que el transistor del buck está encendido. Idealmente
  $D = V_\text{OUT}/V_\text{IN}$.
  \item[Desacoplo] Condensador puesto pegado al pin de alimentación de un integrado. Entrega la corriente rápida que el
  chip pide en cada conmutación, para que no tenga que venir desde lejos por pistas con inductancia.
  \item[ESR] \emph{Equivalent Series Resistance}. Resistencia parásita en serie de un condensador. Limita cuánto filtra
  y cuánto se calienta con corriente de rizado.
  \item[FPWM] \emph{Forced PWM}. Modo en que el buck conmuta siempre a frecuencia fija, aun sin carga. El ruido queda en
  una frecuencia conocida (\SI{1}{\mega\hertz}, lejos del audio), a cambio de un consumo en vacío mayor.
  \item[Lazo de potencia (hot loop)] Camino que recorre la corriente pulsante del buck: condensador de entrada,
  transistor y tierra. Cuanto menor su área, menos radia y menos ruido mete en la placa.
  \item[Rizado] Variación periódica de una tensión o corriente alrededor de su valor medio. En la bobina del buck es el
  triángulo de corriente $\Delta I_L$.
  \item[$t_\text{ON}$ / $t_\text{ON-MIN}$] Tiempo que el transistor del buck está encendido en cada ciclo, y el mínimo
  que el chip puede hacer. Si se necesitara menos, el buck tendría que saltarse pulsos.
  \item[TVS] Diodo de supresión de transitorios. No conduce en operación normal y se enciende cuando la tensión pasa su
  ruptura, absorbiendo el pico. D2 (SMBJ26A) protege la entrada de \SI{24}{\volt}.
  \item[X7R, X5R, C0G, Y5V] Tipos de dieléctrico cerámico. C0G casi no cambia con la tensión ni con la temperatura
  (se usa en los \SI{10}{\pico\farad} de CF). X7R y X5R cambian moderadamente. Y5V pierde gran parte de su capacidad con
  tensión y temperatura; por eso se cambiaron los \SI{100}{\nano\farad} Y5V del diseño original por X7R.
\end{description}

\section*{Amplificadores y conversión}
\begin{description}[style=sameline,leftmargin=1.5em,labelindent=0pt]
  \item[Delta-sigma] Tipo de ADC (como el PCM1808) que muestrea muy rápido con pocos bits y luego filtra digitalmente
  para obtener muchos bits a la frecuencia final. Da alta resolución en audio.
  \item[GBW] \emph{Gain-Bandwidth Product}. Producto ganancia por ancho de banda de un op-amp: con ganancia 1, su ancho
  de banda es más o menos el GBW. El TLV9152 tiene \SI{4.5}{\mega\hertz}.
  \item[Margen de fase] Cuánta fase le falta al lazo de realimentación para volverse oscilador, medida donde la ganancia
  de lazo vale 1. Con menos de unos \SI{45}{\degree} aparecen sobrepicos; cerca de \SI{0}{\degree}, oscila.
  \item[Rail-to-rail] Op-amp cuya entrada (o salida) puede llegar hasta las tensiones de alimentación.
  \item[Región de transición] En un op-amp con dos pares de entrada (P y N), zona de tensión de modo común en la que
  conducen los dos. Ahí empeoran el offset, la CMRR y el ruido. En el TLV9152 va típicamente de $V_+ - \SI{2}{\volt}$
  a $V_+ - \SI{1}{\volt}$.
\end{description}

\section*{Comunicación digital}
\begin{description}[style=sameline,leftmargin=1.5em,labelindent=0pt]
  \item[I2C] Bus de dos hilos (SCL reloj y SDA datos) con resistencias de pull-up. La NUCLEO lo usa para programar los
  DAC MCP4725 (direcciones \texttt{0x60} y \texttt{0x61}).
  \item[I2S / SAI] I2S es el formato de audio digital serie: reloj de bit (BCK), reloj de trama (LRCK, izquierda o
  derecha) y datos (DOUT). SAI es el periférico del STM32 que lo recibe.
  \item[MCLK / SCKI] Reloj maestro del ADC. Aquí lo da el oscilador Y1 de \SI{24.576}{\mega\hertz}, que es
  $512 \times \SI{48}{\kilo\hertz}$.
\end{description}

\section*{La placa: cobre, capas y estructura}
\begin{description}[style=sameline,leftmargin=1.5em,labelindent=0pt]
  \item[Alivio térmico] Unión de un pad a una zona de cobre por medio de unos pocos ``rayos'' en vez de cobre sólido.
  Facilita soldar a mano (el calor no se escapa al plano) pero agrega resistencia e inductancia.
  \item[Anillo anular] Corona de cobre que rodea el agujero de un pad o de una vía. Si es muy delgado, un taladro algo
  corrido lo rompe.
  \item[Apilado (stackup)] Orden y espesor de las capas de cobre y de dieléctrico. Esta placa usa el estándar de JLC para
  4 capas y \SI{1.6}{\milli\meter} (JLC04161H-7628).
  \item[Corriente de retorno] La corriente que vuelve a la fuente por tierra. A frecuencias altas viaja justo debajo de
  la pista de ida, por el plano. Si el plano tiene un corte, la corriente lo rodea y el lazo se agranda.
  \item[Cosido (stitching)] Muchas vías de tierra repartidas que unen el relleno de tierra de una capa con los planos de
  las otras, para que todo sea una sola tierra de baja impedancia.
  \item[Courtyard] Contorno del espacio que ocupa una pieza más un margen. Dos courtyards no deben superponerse.
  \item[Microstrip] Pista en una capa externa sobre un plano de tierra. Su impedancia depende del ancho, del espesor del
  dieléctrico y de $\varepsilon_r$.
  \item[Net class (clase de red)] Grupo de redes con las mismas reglas (ancho, separación, vía). Aquí: Default, Potencia,
  Conmutado, Analogica y Reloj.
  \item[Plano] Capa (o gran parte de una) llena de cobre de una sola red, normalmente tierra. \capa{In1.Cu} es el plano
  de tierra entero de esta placa.
  \item[Stub] Pedazo de pista que no lleva a ningún lado. Es cobre sobrante que puede actuar como antena.
  \item[Teardrop] Ensanchamiento en forma de gota donde una pista llega a un pad o a una vía. Refuerza la unión
  mecánica y tolera mejor un taladro corrido.
  \item[Vía] Agujero metalizado que une capas. En esta placa: \SI{0.3}{\milli\meter} de taladro y \SI{0.6}{\milli\meter}
  de pad.
  \item[Vía taponada (plugged) / tenting] Vía cubierta: \emph{tenting} es taparla con máscara; \emph{plugged} es rellenar
  el agujero antes. JLC tapona las vías de 4 capas, y por eso no pueden quedar a menos de \SI{0.35}{\milli\meter} de
  un pad.
  \item[Zona] Área de cobre que KiCad rellena automáticamente respetando las separaciones. Tiene prioridad: la de mayor
  prioridad se rellena primero y las demás la rodean.
\end{description}

\section*{Fabricación y montaje}
\begin{description}[style=sameline,leftmargin=1.5em,labelindent=0pt]
  \item[Basic / Extended] Clasificación de piezas de JLCPCB. Las \emph{basic} están siempre cargadas en las máquinas;
  cada \emph{extended} distinta cobra una tarifa de carga del alimentador.
  \item[BOM] \emph{Bill of Materials}. Lista de materiales: qué pieza va en cada designador y su código LCSC.
  \item[CPL] \emph{Component Placement List}. Posición, cara y rotación de cada pieza para la máquina de montaje.
  \item[DFM] \emph{Design for Manufacturing}. Revisión de que la placa se pueda fabricar y montar; JLC muestra una vista
  previa DFM antes de pagar.
  \item[Economic / Standard PCBA] Dos servicios de montaje de JLC. Economic es más barato, solo monta una cara y junta
  pedidos de varios clientes. Standard monta las dos caras y pide placas de al menos \SI{70}{} $\times$
  \SI{70}{\milli\meter}.
  \item[ENIG] Acabado de oro sobre níquel. Muy plano y caro.
  \item[Excellon] Formato de los archivos de taladro (posición y diámetro de cada agujero).
  \item[Fiducial] Marca de cobre redonda sin máscara alrededor que la cámara de la máquina de montaje usa para
  ubicar la placa. Esta tiene tres.
  \item[Footprint] Dibujo de los pads, la serigrafía y el contorno de una pieza en la placa. Debe coincidir con el
  encapsulado real y con el que usa el ensamblador.
  \item[Gerber] Formato estándar de los archivos de fabricación, uno por capa (cobre, máscara, serigrafía, pasta,
  borde).
  \item[HASL] Acabado de estaño fundido nivelado con aire caliente. Barato y fácil de soldar; existe con plomo y sin
  plomo. Esta placa lleva HASL sin plomo.
  \item[LCSC] Distribuidor de componentes asociado a JLCPCB. El código LCSC (por ejemplo C2861011) identifica la pieza
  exacta que monta JLC.
  \item[Máscara de soldadura] Capa de laca (verde en esta placa) que cubre el cobre salvo los pads. Evita puentes de
  soldadura y protege el cobre.
  \item[Puente de máscara] Franja de máscara entre dos pads vecinos. Si es más angosta que el mínimo del fabricante,
  desaparece y la soldadura puede unir los pads.
  \item[Riel (edge rail)] Franja de material que se agrega al borde de la placa para que la máquina de montaje la sujete.
  Se corta después. Aquí los agrega JLC.
  \item[Serigrafía] Capa de tinta (blanca) con designadores, contornos y marcas de polaridad. Nunca debe quedar sobre
  un pad.
\end{description}

\section*{Verificación}
\begin{description}[style=sameline,leftmargin=1.5em,labelindent=0pt]
  \item[DRC] \emph{Design Rule Check}. Revisión automática de la geometría de la placa contra las reglas (separaciones,
  anchos, anillos, bordes).
  \item[ERC] \emph{Electrical Rule Check}. Revisión del esquemático según el tipo de cada pin (salidas unidas, entradas
  al aire, alimentaciones sin fuente).
  \item[Paridad] Comparación entre la placa y el esquemático: mismas piezas, mismos footprints y mismas conexiones.
  \item[Ratsnest] Líneas rectas que KiCad dibuja entre pads que deben unirse y todavía no tienen pista.
\end{description}
